\chapter{Metodologia}
\label{chap:metodologia}

Para atingir os objetivos propostos neste trabalho, notadamente a investigação das percepções e dos fatores que influenciam a adoção de JavaScript e TypeScript entre diferentes perfis de desenvolvedores, optou-se por uma abordagem de métodos mistos, com predominância quantitativa e complemento qualitativo. O estudo é caracterizado como uma pesquisa aplicada, de natureza descritiva e comparativa, realizada por meio de levantamento (\textit{survey}). O desenho quantitativo permite sintetizar respostas padronizadas, enquanto a análise qualitativa dos relatos abertos ajuda a explicar por que determinadas percepções aparecem nos dados. Essa combinação é coerente com o problema, pois a preferência tecnológica possui uma dimensão mensurável e outra relacionada às experiências de trabalho, manutenção e colaboração.

\section{Problema de pesquisa e questões investigativas}

O problema de pesquisa pode ser formulado da seguinte maneira: como desenvolvedores com diferentes níveis de experiência percebem os benefícios, custos e riscos associados ao uso de JavaScript e TypeScript no desenvolvimento front-end? A pergunta não pressupõe que uma tecnologia seja universalmente superior. Ela considera que prototipagem, aprendizagem, manutenção, refatoração, integração e governança podem produzir avaliações diferentes.

Com base nesse problema, o levantamento foi orientado pelas seguintes questões investigativas:
\begin{enumerate}
	\item Como o grupo iniciante/intermediário e o grupo experiente avaliam produtividade, facilidade de aprendizagem e rapidez de prototipagem?
	\item Como os grupos avaliam detecção precoce de erros, legibilidade, manutenibilidade e padronização?
	\item Que dificuldades são percebidas na configuração, migração e integração com bibliotecas externas?
	\item Em que medida a experiência em projetos de maior escala influencia a preferência pelo TypeScript?
	\item Que benefícios e limitações aparecem nos relatos abertos sobre uso real das duas tecnologias?
\end{enumerate}

Essas questões dialogam com a literatura sobre dívida técnica. Uma decisão que reduz o esforço inicial pode produzir custos posteriores de manutenção, enquanto uma configuração mais rigorosa pode exigir investimento antecipado para reduzir problemas futuros \cite{alves2016debt}. Também dialogam com estudos que mostram que a percepção de gestores e desenvolvedores pode enfatizar dimensões diferentes da dívida técnica \cite{gomes2023debt}.

\section{Instrumento de coleta}

O instrumento de coleta de dados consistiu em dois questionários distintos, elaborados e distribuídos por meio da plataforma Google Forms. A decisão de criar dois formulários separados foi estratégica, visando adequar a linguagem e o foco das perguntas ao nível de conhecimento e ao contexto de cada público-alvo, sem alterar o núcleo das dimensões comparadas. Ambos os questionários foram estruturados com perguntas fechadas, utilizando escalas de Likert de cinco pontos, perguntas de múltipla escolha e uma questão aberta sobre experiências práticas.

As afirmações foram organizadas em quatro blocos analíticos. O Bloco A avaliou produtividade, agilidade, curva de aprendizagem, prototipagem e detecção precoce de erros. O Bloco B avaliou legibilidade, manutenibilidade, origem de bugs, verbosidade e padronização. O Bloco C abordou compatibilidade com bibliotecas, ferramentas, frameworks, migração e projetos legados. O Bloco D examinou confiança, preferência para novos projetos, percepção de futuro e adequação a equipes grandes. A pergunta aberta foi incluída para identificar situações não previstas nas alternativas fechadas.

Essa estrutura procura operacionalizar os principais construtos do referencial. A detecção precoce representa o benefício esperado da verificação estática; manutenibilidade e refatoração se relacionam ao custo de evolução; migração e configuração representam o investimento inicial; e integração com bibliotecas representa a confiabilidade dos contratos externos. A inclusão de processo, equipe e colaboração é importante porque a literatura mostra que a dívida técnica também envolve pessoas e decisões organizacionais, não apenas o código \cite{gomes2023debt,coutinho2024agile}.

\subsection{Equivalência entre os formulários}

Embora tenham sido utilizados dois formulários, os instrumentos mantiveram as mesmas dimensões, itens e escalas para permitir a comparação entre os grupos. As diferenças entre as versões limitaram-se à apresentação das instruções e à adequação de exemplos ao perfil do respondente. Assim, os itens dos blocos A, B, C e D foram tratados como correspondentes nos dois questionários. Essa equivalência é necessária para que diferenças nas médias sejam interpretadas como diferenças de percepção, e não como efeito de perguntas distintas.

\subsection{Construção e verificação do instrumento}

Os itens foram elaborados a partir dos temas recorrentes no referencial teórico e nos objetivos da pesquisa: tipagem, detecção de erros, produtividade percebida, manutenibilidade, migração, integração com bibliotecas, ferramentas e colaboração. A organização em blocos foi revisada para reduzir sobreposição entre dimensões e facilitar a interpretação dos resultados. Não foi realizada validação psicométrica formal nem estimativa de confiabilidade interna; por isso, os resultados são tratados como descrição exploratória das respostas ao instrumento, e não como uma escala validada para medir construtos em qualquer população.
\section{População, amostragem e grupos de análise}

O primeiro grupo-alvo foi composto por \textbf{desenvolvedores profissionais experientes}, atuantes no mercado de desenvolvimento de software. A amostra foi selecionada por conveniência, utilizando redes de contatos profissionais, como o LinkedIn, e comunidades online de desenvolvedores. Para manter a coerência com a segmentação apresentada nos resultados, foram considerados experientes, exclusivamente como categoria operacional desta pesquisa, os participantes com quatro ou mais anos de experiência declarada em desenvolvimento front-end. O grupo é relevante para compreender percepções formadas em projetos comerciais, de longa duração e com equipes colaborativas, nos quais manutenibilidade, escalabilidade e governança assumem maior importância.

O segundo grupo foi formado por \textbf{desenvolvedores iniciantes e intermediários}, majoritariamente estudantes do curso de Bacharelado em Sistemas de Informação do Instituto Federal da Bahia (IFBA), Campus Feira de Santana, e de outras instituições. Para a comparação principal, foram agrupados participantes com até três anos de experiência. A seleção também foi por conveniência, por meio de convites em grupos de estudo e canais acadêmicos. Esse grupo representa participantes ainda próximos do processo de aprendizagem e, por isso, é especialmente adequado para investigar curva de aprendizado, simplicidade inicial e percepção do custo de configuração \cite{shrestha2022}.

A amostra final foi composta por 23 respondentes: 14 no grupo iniciante/intermediário e 9 no grupo experiente. Doze participantes declararam experiência em projetos de grande escala e 11 relataram atuação predominantemente em projetos individuais, acadêmicos ou de menor porte. Essa variável foi utilizada como dimensão complementar, pois a escala do sistema altera a relevância de contratos, refatorações e coordenação entre pessoas.

\section{Procedimentos de coleta e tratamento dos dados}

Após a divulgação dos formulários, as respostas foram exportadas para planilhas eletrônicas e submetidas a conferência. Foram verificadas duplicidades, respostas incompletas e a correspondência entre as respostas e o grupo de experiência informado. Os dados foram então codificados numericamente, preservando a direção da escala: valores maiores representam maior concordância com a afirmação apresentada. As respostas abertas foram mantidas em texto para a etapa qualitativa.

Não se pretendeu produzir uma estimativa probabilística da população de desenvolvedores. A amostragem por conveniência permite explorar tendências e contrastes dentro do grupo alcançado, mas exige cautela na generalização. Por esse motivo, os resultados são interpretados como evidências descritivas sobre as percepções dos respondentes, e não como prova de que todos os participantes com menor ou maior experiência se comportam da mesma maneira.

\section{Plano de análise quantitativa}

A análise quantitativa utilizou estatística descritiva, incluindo frequências absolutas e relativas, médias, medianas e desvios-padrão. Para cada item, foram calculadas as respostas da amostra total e dos dois grupos principais. A diferença perceptual entre experientes e iniciantes foi expressa por:
\[
\Delta = \overline{x}_{experientes} - \overline{x}_{iniciantes}.
\]

Valores positivos indicam maior concordância entre experientes; valores negativos indicam maior concordância entre iniciantes. A mediana foi apresentada junto à média porque a escala de Likert é ordinal e pode concentrar respostas em determinados pontos. O desvio-padrão foi usado como indicador descritivo de dispersão, permitindo distinguir consenso de respostas mais heterogêneas.

As comparações foram organizadas por dimensão, e não apenas por item isolado. Essa decisão reduz o risco de interpretar uma resposta específica fora do contexto e permite relacionar os dados aos construtos teóricos: produtividade, qualidade, integração e preferência. As diferenças foram discutidas em conjunto com o tamanho dos grupos, a experiência em projetos grandes e os relatos qualitativos. Não foram feitas afirmações causais a partir das médias, pois o desenho transversal e a amostragem não permitem concluir que a experiência provoque determinada preferência.

Não foram realizados testes de hipótese ou inferência estatística. Essa decisão decorre do caráter descritivo e exploratório do estudo, do uso de amostragem por conveniência e do tamanho reduzido dos grupos, com 14 participantes no grupo iniciante/intermediário e 9 no grupo experiente. Portanto, o valor de $\Delta$ é apresentado como diferença descritiva entre as respostas observadas, e não como evidência de diferença estatisticamente significativa na população.

\section{Plano de análise qualitativa}

As respostas abertas foram examinadas por análise temática categorial. Inicialmente, foram identificadas unidades de sentido relacionadas a benefícios, custos, contexto de uso, integração, manutenção e execução. Em seguida, os relatos foram agrupados em quatro eixos: blindagem de contratos e APIs; refatoração em larga escala; complexidade inicial e fricção com o ecossistema; e fronteira entre compilação e execução.

Os eixos não foram tratados como categorias mutuamente exclusivas de defeitos. Um mesmo relato poderia mencionar tipagem, biblioteca externa e validação de runtime. A categorização teve como finalidade aproximar a experiência dos participantes dos conceitos encontrados na literatura, especialmente dívida técnica, contratos quebrados, toolchain e colaboração. A interpretação buscou preservar o sentido dos depoimentos e evitar transformar opiniões individuais em generalizações sobre toda a comunidade.

\section{Validade, limitações e aspectos éticos}

Há ameaças à validade interna, externa e de construção. A validade interna é limitada porque as respostas são autorreferidas e podem refletir preferência pessoal, familiaridade com ferramentas ou experiências específicas. A validade externa é limitada pela amostragem por conveniência e pelo número reduzido de participantes. A validade de construção exige cautela porque termos como produtividade, segurança de tipos e manutenibilidade são complexos e foram medidos por percepção, não por métricas de repositórios ou experimentos controlados.

Para reduzir esses riscos, foram utilizadas perguntas organizadas por dimensões, comparação entre grupos, médias acompanhadas de medianas e desvios-padrão e integração dos resultados fechados com os relatos abertos. A literatura também foi usada para interpretar os resultados sem reduzir qualidade a uma única variável. Estudos de mineração de repositórios, como os que investigam opções do compilador e categorias de bugs, servem como referência externa, mas não substituem a análise do survey, pois empregam unidades de análise e objetivos diferentes.

No plano ético, a participação foi voluntária e os dados foram analisados de forma agregada. Os relatos apresentados no capítulo de resultados são identificados por códigos de respondentes, sem exposição de nomes ou informações pessoais desnecessárias. O uso de depoimentos teve finalidade acadêmica e foi restrito ao conteúdo relacionado ao problema de pesquisa.

Além do tamanho reduzido, os grupos apresentam distribuição desigual, o que recomenda cautela na comparação direta. A familiaridade prévia com uma das linguagens pode influenciar as respostas e produzir viés de preferência. Também não foram coletadas métricas objetivas de tempo de desenvolvimento, densidade de bugs, complexidade, rotatividade de código ou custo de migração. Assim, afirmações sobre produtividade, qualidade e manutenção representam percepções declaradas, e não medições experimentais ou observacionais do desempenho dos projetos.

\section{Síntese metodológica}

Em síntese, o método combina levantamento quantitativo, comparação entre níveis de experiência e interpretação qualitativa de relatos. Essa estratégia é adequada ao objetivo de investigar não apenas qual linguagem recebe maior preferência, mas quais razões sustentam essa preferência. A abordagem também permite compreender a escolha como decisão sociotécnica: a linguagem interage com configuração, contratos, arquitetura, processo, aprendizagem e colaboração. Dessa forma, os resultados podem ser confrontados com a literatura de qualidade e dívida técnica sem extrapolar o alcance da amostra.

